\documentclass[11pt]{article}
\usepackage[a4paper,margin=0.9in]{geometry}
\usepackage{amsmath,amssymb,amsthm,mathtools}
\usepackage[dvipsnames]{xcolor}
\usepackage{enumitem}
\usepackage{booktabs}
\usepackage{longtable}
\usepackage{tcolorbox}
\usepackage{fancyhdr}
\usepackage[colorlinks=true,linkcolor=Blue]{hyperref}
\tcbuselibrary{skins,breakable}

\setlist{itemsep=3pt,topsep=4pt}
\pagestyle{fancy}\fancyhf{}
\lhead{\small Minimax \& Theorem Compendium}
\rhead{\small Parametric Inference, B.Stat.\ 3rd yr}
\cfoot{\small\thepage}

\newcommand{\E}{\mathbb{E}}
\newcommand{\Var}{\operatorname{Var}}
\newcommand{\Cov}{\operatorname{Cov}}
\newcommand{\Prob}{\mathbb{P}}
\newcommand{\R}{\mathbb{R}}
\newcommand{\Xbar}{\overline{X}}
\newcommand{\iid}{\overset{\text{iid}}{\sim}}
\newcommand{\pto}{\overset{p}{\longrightarrow}}
\newcommand{\ind}{\perp\!\!\!\perp}
\newcommand{\I}{\mathbb{I}}

\newtcolorbox{idea}[1]{colback=Blue!4,colframe=Blue!55,title={#1},
  fonttitle=\bfseries,breakable}
\newtcolorbox{confusion}[1]{colback=BrickRed!5,colframe=BrickRed!60,
  title={Common confusion: #1},fonttitle=\bfseries,breakable}
\newtcolorbox{recipe}[1]{colback=ForestGreen!5,colframe=ForestGreen!55,title={#1},
  fonttitle=\bfseries,breakable}

\newcounter{thm}
\newenvironment{Thm}[1]{\medskip\refstepcounter{thm}%
  \noindent\textbf{\large T\thethm.\ #1}\par\nopagebreak\smallskip}{\par\medskip}
\newcommand{\hyp}{\smallskip\noindent\textbf{\color{BrickRed}Hypotheses (state and check these):}\ }
\newcommand{\conc}{\smallskip\noindent\textbf{Conclusion:}\ }
\newcommand{\prf}{\smallskip\noindent\textbf{Proof:}\ }
\newcommand{\uses}{\smallskip\noindent\textbf{\color{ForestGreen}Used for:}\ }

\begin{document}

\begin{center}
{\LARGE\bfseries Minimax Estimation}\\[3pt]
{\large A Crash Course --- plus a Compendium of Every Theorem}\\[5pt]
{\small Parametric Inference, B.Stat.\ 3rd Year --- Lectures 8--9, and all ten lectures}
\end{center}

\vspace{4pt}\hrule\vspace{10pt}

\noindent
Third and last of the crash courses, after \texttt{PI\_Bayes\_Crash\_Course.pdf} and
\texttt{PI\_Testing\_Crash\_Course.pdf}. \textbf{Part I} builds minimax from scratch;
\textbf{Part II} is a reference containing every theorem in the course, each with its
hypotheses called out separately --- because in this subject the marks are in the hypotheses,
not the statement.

\tableofcontents

\newpage
\part*{Part I --- Minimax from scratch}
\addcontentsline{toc}{part}{Part I --- Minimax from scratch}

%%%%%%%%%%%%%%%%%%%%%%%%%%%%%%%%%%%%%%%%%%%%%%%%%%%%%%%%%
\section{The third escape route}

The same starting point as always. The risk $R(\theta,T)=\E_\theta(T-\theta)^2$ is a
\emph{function} of $\theta$, so ``minimise the risk'' is not a well-posed problem. Three ways
to make it well posed, and you now know all three:

\begin{center}
\renewcommand{\arraystretch}{1.4}
\begin{tabular}{@{}p{2.6cm}p{5.6cm}p{7.3cm}@{}}
\toprule
& \textbf{How the function becomes a number} & \textbf{Character} \\
\midrule
\textbf{UMVUE} & Don't collapse it --- shrink the class of $T$ instead & Optimistic: assumes unbiasedness is a reasonable restriction \\
\textbf{Bayes} & $r(\pi,T)=\int R(\theta,T)\pi(\theta)d\theta$ --- \emph{average} it & Requires you to commit to weights $\pi$ \\
\textbf{Minimax} & $\sup_\theta R(\theta,T)$ --- take the \emph{worst case} & Pessimistic: protects against the least favourable $\theta$, commits to nothing \\
\bottomrule
\end{tabular}
\end{center}

\begin{idea}{Definition}
$T^*$ is \textbf{minimax} if
\[
\sup_{\theta\in\Theta}R(\theta,T^*)\ =\ \inf_{T}\ \sup_{\theta\in\Theta}R(\theta,T).
\]
Read it right to left: among all estimators, look at each one's worst-case risk, and pick the
estimator whose worst case is least bad. The common value is the \textbf{minimax value} of the
problem.
\end{idea}

%%%%%%%%%%%%%%%%%%%%%%%%%%%%%%%%%%%%%%%%%%%%%%%%%%%%%%%%%
\section{Why minimax is hard, and the one trick that solves it}

Suppose someone hands you a candidate $T^*$. Computing $\sup_\theta R(\theta,T^*)$ is easy ---
it is calculus. But proving $T^*$ is \emph{minimax} means showing
\[
\sup_\theta R(\theta,T)\ \ge\ \sup_\theta R(\theta,T^*)\qquad\textbf{for every estimator }T,
\]
and there are infinitely many $T$, most of which you cannot write down. You cannot check them
one at a time. You need a \textbf{certificate}: a lower bound on $\inf_T\sup_\theta R(\theta,T)$
that you can compute, and that happens to equal $\sup_\theta R(\theta,T^*)$.

\begin{idea}{Where lower bounds come from --- this is the whole subject}
For \textbf{any} probability measure $\pi$ on $\Theta$ and \textbf{any} estimator $T$,
\[
\sup_{\theta}R(\theta,T)\ \ge\ \int_\Theta R(\theta,T)\,\pi(\theta)\,d\theta\ =\ r(\pi,T),
\]
simply because a supremum is at least any weighted average. Now minimise both sides over $T$:
\[
\inf_T\sup_\theta R(\theta,T)\ \ge\ \inf_T r(\pi,T)\ =\ r(\pi,T_\pi),
\]
the Bayes risk of the Bayes estimate for $\pi$.

\smallskip
\textbf{So every prior $\pi$ hands you a free lower bound on the minimax value, namely its
Bayes risk.} That is why Bayes estimation is taught immediately before minimax --- it is not a
detour, it is the machine that produces the certificates.
\end{idea}

\begin{idea}{The sandwich}
Putting it together, for any $\pi$ and any $T$:
\[
\underbrace{\sup_\theta R(\theta,T^*)}_{\text{what you can compute}}
\ \ \ge\ \ \underbrace{\inf_T\sup_\theta R(\theta,T)}_{\text{what you want}}
\ \ \ge\ \ \underbrace{r(\pi,T_\pi)}_{\text{a computable lower bound}} .
\]
The left inequality is trivial ($T^*$ is one of the $T$'s). \textbf{If you can find a $\pi$
making the two outer quantities equal, the sandwich collapses and $T^*$ is minimax.}
Everything in Lectures 8--9 is a recipe for finding that $\pi$.

A $\pi$ achieving the best possible lower bound is called a \textbf{least favourable prior}:
it puts its mass where estimation is hardest, making the problem as difficult as it can be.
\end{idea}

%%%%%%%%%%%%%%%%%%%%%%%%%%%%%%%%%%%%%%%%%%%%%%%%%%%%%%%%%
\section{The two theorems, and when to use which}

\begin{idea}{Theorem A --- constant-risk (equalizer) criterion}
If $T_\pi$ is the Bayes estimate for a \textbf{proper} prior $\pi$, and
$R(\theta,T_\pi)=c$ is \textbf{constant in $\theta$}, then $T_\pi$ is minimax.
\end{idea}

\prf Substitute into the sandwich. For any $T$,
\[
\sup_\theta R(\theta,T)\ \ge\ r(\pi,T)\ \ge\ r(\pi,T_\pi)
=\int c\,\pi(\theta)d\theta=c=\sup_\theta R(\theta,T_\pi).\qquad\square
\]
The constant risk is what makes step three work: the average of a constant is that constant,
so the lower bound lands exactly on $T_\pi$'s worst case.

\begin{idea}{Theorem B --- limit of Bayes rules}
Let $\{\pi_k\}$ be proper priors with Bayes risks $r(\pi_k,T_{\pi_k})\to c$. If $T^*$
satisfies $\sup_\theta R(\theta,T^*)=c$, then $T^*$ is minimax.
\end{idea}

\prf For any $T$ and every $k$, $\sup_\theta R(\theta,T)\ge r(\pi_k,T)\ge r(\pi_k,T_{\pi_k})$.
Let $k\to\infty$: $\sup_\theta R(\theta,T)\ge c=\sup_\theta R(\theta,T^*)$. $\square$

\begin{confusion}{``Why do we need Theorem B at all?''}
Because \textbf{the minimax estimator is often not Bayes for any proper prior}, so Theorem A
simply does not apply.

The example is $X\sim N(\theta,1)$ with $\theta\in\R$. The minimax estimator is $X$ itself,
with constant risk 1. But is $X$ Bayes for some proper prior? No: every proper prior pulls the
estimate toward its own mean, so every proper Bayes estimate is a strict shrinkage
$\frac{k}{k+1}X$, never $X$. Only the \emph{improper} flat prior gives $X$, and an improper
prior has infinite Bayes risk, so the ``$\sup\ \ge$ average'' step is meaningless.

The repair is Theorem B: use proper priors $\pi_k=N(0,k)$ that approach flatness, and let
their Bayes risks $k/(k+1)$ climb to $1$. The sequence $\{\pi_k\}$ is called a \emph{least
favourable sequence}. \textbf{Rule of thumb: bounded $\Theta$ $\to$ try Theorem A; unbounded
$\Theta$ $\to$ you will almost certainly need Theorem B.}
\end{confusion}

\begin{idea}{Why minimax rules tend to have constant risk}
Not a theorem, but the right intuition. Suppose $R(\cdot,T)$ is not flat, peaking at some
$\theta_1$. Intuitively you can nudge $T$ to do a little worse where the risk is low and a
little better near $\theta_1$, lowering the peak --- and the peak is all that minimaxity cares
about. At the optimum there should be nothing left to shave, i.e.\ the risk is flat.

An estimator with constant risk is called an \textbf{equalizer}. So the practical search
strategy is: \emph{look for an equalizer inside a conjugate family, then check it is Bayes.}
That is exactly how the binomial answer is found below.
\end{idea}

%%%%%%%%%%%%%%%%%%%%%%%%%%%%%%%%%%%%%%%%%%%%%%%%%%%%%%%%%
\section{Worked example 1: the Binomial (Theorem A in full)}

$X\sim\mathrm{Bin}(n,\theta)$, $\theta\in(0,1)$. Here $\Theta$ is bounded, so try Theorem A.

\textbf{Step 1. Get a family of Bayes estimates.} With prior $\mathrm{Beta}(p,q)$ the posterior
is $\mathrm{Beta}(p+x,\ q+n-x)$, so writing $c=p+q+n$,
\[
T_{p,q}=\frac{p+X}{c}.
\]

\textbf{Step 2. Compute its risk.} $\E_\theta X=n\theta$ and $\Var_\theta X=n\theta(1-\theta)$,
so
\[
\text{Bias}=\frac{p+n\theta}{c}-\theta=\frac{p-(p+q)\theta}{c},
\qquad
\Var(T_{p,q})=\frac{n\theta(1-\theta)}{c^{2}},
\]
\[
R(\theta,T_{p,q})=\frac{1}{c^{2}}\Big[n\theta(1-\theta)+\big(p-(p+q)\theta\big)^{2}\Big].
\]

\textbf{Step 3. Force it to be constant.} Expand the bracket as a quadratic in $\theta$:
\[
\underbrace{\big[(p+q)^{2}-n\big]}_{\text{coeff of }\theta^{2}}\theta^{2}
+\underbrace{\big[n-2p(p+q)\big]}_{\text{coeff of }\theta}\theta
+\underbrace{p^{2}}_{\text{constant}} .
\]
Both non-constant coefficients must vanish:
\[
(p+q)^{2}=n\ \Rightarrow\ p+q=\sqrt n;
\qquad
n=2p(p+q)=2p\sqrt n\ \Rightarrow\ p=\frac{\sqrt n}{2},\ q=\frac{\sqrt n}{2}.
\]
This is the \emph{only} solution --- two equations, two unknowns.

\textbf{Step 4. Read off the answer.} With $c=n+\sqrt n=\sqrt n(1+\sqrt n)$,
\[
\boxed{\ T^*=\frac{X+\sqrt n/2}{n+\sqrt n},\qquad
R(\theta,T^*)=\frac{p^{2}}{c^{2}}=\frac{n/4}{n(1+\sqrt n)^{2}}=\frac{1}{4(1+\sqrt n)^{2}}\ }
\]
for every $\theta$. It is Bayes for a proper prior and has constant risk, so by Theorem A it is
\textbf{minimax}, and $\mathrm{Beta}(\sqrt n/2,\sqrt n/2)$ is least favourable.

\begin{idea}{Sanity check on the least favourable prior}
$\mathrm{Beta}(\sqrt n/2,\sqrt n/2)$ has mean $\tfrac12$ and variance
$\dfrac{n/4}{n(\sqrt n+1)}=\dfrac{1}{4(\sqrt n+1)}\to0$. So as $n$ grows the least favourable
prior \textbf{concentrates at $\theta=\tfrac12$} --- exactly where $\Var_\theta(X)=n\theta(1-\theta)$
is largest, i.e.\ where estimation is hardest. The certificate is being placed precisely where
the problem is worst, which is what ``least favourable'' should mean.
\end{idea}

\subsection{Minimax versus the UMVUE --- and what minimaxity actually buys}

The UMVUE is $X/n$, with $R(\theta,X/n)=\theta(1-\theta)/n$ and
$\max_\theta R=\dfrac{1}{4n}$ at $\theta=\tfrac12$. Since
$(1+\sqrt n)^2=n+2\sqrt n+1>n$,
\[
\frac{1}{4(1+\sqrt n)^{2}}\ <\ \frac{1}{4n}
\quad\Longrightarrow\quad \textbf{the UMVUE is not minimax.}
\]
But look at the whole risk curves, say $n=100$:
\begin{center}
\renewcommand{\arraystretch}{1.3}
\begin{tabular}{@{}lccc@{}}
\toprule
& $\theta=0.5$ & $\theta=0.1$ & $\theta\to0$ \\
\midrule
UMVUE $X/n$ & $0.00250$ & $0.00090$ & $\to0$ \\
Minimax $T^*$ & $0.002066$ & $0.002066$ & $0.002066$ \\
\bottomrule
\end{tabular}
\end{center}
$T^*$ wins at the peak and \textbf{loses badly everywhere else} --- near the endpoints the
UMVUE has almost no risk while $T^*$ still pays its flat premium.

\begin{confusion}{``So is the minimax estimator the best one?''}
No. Minimaxity is a \textbf{worst-case} guarantee, not a uniform one. $T^*$ is the right choice
only if you genuinely fear an adversary choosing the worst $\theta$. If you have any reason to
think $\theta$ is near $0$ or $1$, the UMVUE is far better. Note also the ratio
$\frac{(1+\sqrt n)^2}{n}\to1$: for large $n$ the two are almost identical anyway, so the whole
distinction is a small-sample phenomenon.
\end{confusion}

%%%%%%%%%%%%%%%%%%%%%%%%%%%%%%%%%%%%%%%%%%%%%%%%%%%%%%%%%
\section{Worked example 2: the Normal (Theorem B in full)}

$X\sim N(\theta,1)$, $\theta\in\R$ --- now $\Theta$ is unbounded.

\textbf{Step 1. Candidate and its worst-case risk.} $R(\theta,X)=\E_\theta(X-\theta)^2=1$ for
every $\theta$, so $X$ is an equalizer with $\sup_\theta R(\theta,X)=1$.

\textbf{Step 2. Build the certificate sequence.} Take $\pi_k=N(0,k)$. Completing the square,
\[
\theta\mid X=x\ \sim\ N\Big(\frac{k}{k+1}x,\ \frac{k}{k+1}\Big),
\qquad T_{\pi_k}=\frac{k}{k+1}X .
\]
The posterior variance is the \emph{constant} $\frac{k}{k+1}$, so by
$r(\pi,T_\pi)=\E_m[\Var(\theta\mid X)]$,
\[
r(\pi_k,T_{\pi_k})=\frac{k}{k+1}.
\]

\textbf{Step 3. Close the sandwich.} $r(\pi_k,T_{\pi_k})\to1=\sup_\theta R(\theta,X)$, so by
Theorem B, $X$ is \textbf{minimax} and the minimax value is $1$.

\begin{idea}{What the sequence is doing}
$\pi_k=N(0,k)$ spreads out as $k$ grows, approaching the (improper) flat prior. Each $\pi_k$
gives a valid lower bound $k/(k+1)$; none of them reaches $1$, but together they climb to it.
No \emph{single} prior certifies $X$; the sequence does. That is the entire reason Theorem B
exists, and it is Homework 2 Q1.
\end{idea}

%%%%%%%%%%%%%%%%%%%%%%%%%%%%%%%%%%%%%%%%%%%%%%%%%%%%%%%%%
\section{Worked example 3: when the minimax value is infinite}

$X\sim\mathrm{Poisson}(\lambda)$, $\lambda\in(0,\infty)$. Take $\pi=\mathrm{Gamma}(a,\text{scale }b)$.
The posterior is $\mathrm{Gamma}\big(a+x,\tfrac{b}{1+b}\big)$ with variance
$\frac{(a+x)b^2}{(1+b)^2}$, and the marginal mean of $X$ is $\E[\lambda]=ab$, so
\[
r(\pi_{a,b})=\E_m\big[\Var(\lambda\mid X)\big]=\frac{b^{2}(a+ab)}{(1+b)^{2}}=\frac{ab^{2}}{1+b}.
\]
Take $a=1$ and $b=k\to\infty$: $r=\dfrac{k^{2}}{1+k}\to\infty$. Since every estimator's
worst-case risk dominates every Bayes risk,
\[
\sup_{\lambda>0}R(\lambda,T)\ \ge\ \frac{k^{2}}{1+k}\ \longrightarrow\ \infty
\qquad\textbf{for every }T .
\]
So \textbf{every} estimator has infinite maximum risk and no minimax estimator exists in any
useful sense (Homework 2 Q2).

\begin{idea}{The pattern to recognise}
Infinite minimax value happens when $\Theta$ is \textbf{unbounded} \emph{and} the achievable
variance grows without bound along it. Two instances in the course:
\begin{itemize}
\item $\mathrm{Poisson}(\lambda)$, $\lambda\in(0,\infty)$: $\Var_\lambda(X)=\lambda\to\infty$.
\item $U[a,\theta]$, $\theta\ge a$ unbounded: with Pareto priors of scale $m$, the Bayes risk
is $\ge C(\alpha,n)m^{2}\to\infty$.
\end{itemize}
Contrast $\mathrm{Bin}(n,\theta)$ with $\theta\in(0,1)$ bounded, and $N(\theta,1)$ where
$\Theta$ is unbounded but the variance is the \emph{constant} 1. \textbf{Unboundedness alone is
not fatal --- it has to be paired with unbounded variance.}
\end{idea}

%%%%%%%%%%%%%%%%%%%%%%%%%%%%%%%%%%%%%%%%%%%%%%%%%%%%%%%%%
\section{The recipe}

\begin{recipe}{Given a minimax question, do this}
\begin{enumerate}
\item \textbf{Is $\Theta$ bounded?} If not, first ask whether the variance blows up along it
--- if so, show every estimator has infinite maximum risk (Poisson pattern) and stop.
\item \textbf{Compute $\sup_\theta R(\theta,T^*)$} for your candidate. If the risk is constant,
you have an equalizer and are in good shape.
\item \textbf{Is $T^*$ Bayes for a proper prior?} Search inside the conjugate family: write
down $R(\theta,T_{\text{prior params}})$, expand as a polynomial in $\theta$, and force the
non-constant coefficients to vanish. If a proper prior works, quote \textbf{Theorem A} and
finish in one line.
\item \textbf{If no proper prior gives $T^*$} (typical when $T^*$ is generalized Bayes under an
improper prior), build a sequence $\pi_k$ of proper priors with $r(\pi_k)\to\sup_\theta
R(\theta,T^*)$ and quote \textbf{Theorem B}.
\item \textbf{Say which prior (or sequence) is least favourable.} It is usually worth a mark,
and it shows you understand what the certificate is.
\end{enumerate}
\end{recipe}

\begin{center}
\renewcommand{\arraystretch}{1.45}
\begin{tabular}{@{}p{3.0cm}p{2.4cm}p{3.5cm}p{6.2cm}@{}}
\toprule
Problem & Minimax rule & Minimax value & Certificate \\
\midrule
$\mathrm{Bin}(n,\theta)$ & $\dfrac{X+\sqrt n/2}{n+\sqrt n}$ & $\dfrac{1}{4(1+\sqrt n)^2}$ &
Theorem A with $\mathrm{Beta}(\tfrac{\sqrt n}2,\tfrac{\sqrt n}2)$ \\
$N(\theta,1)$ & $X$ & $1$ & Theorem B with $\pi_k=N(0,k)$ \\
$N(\theta,\sigma_0^2)$, $n$ obs & $\Xbar$ & $\sigma_0^2/n$ & Theorem B, same idea \\
$\mathrm{Poi}(\lambda)$ & none & $\infty$ & Gamma priors, $r\to\infty$ \\
$U[a,\theta]$ & none & $\infty$ & Pareto priors, $r\to\infty$ \\
\bottomrule
\end{tabular}
\end{center}

\newpage
\part*{Part II --- Every theorem in the course}
\addcontentsline{toc}{part}{Part II --- Every theorem in the course}

\section{Quick reference}

The hypotheses column is what examiners actually check. Quoting a theorem without verifying
its hypotheses is the most reliable way to lose marks in this subject.

\renewcommand{\arraystretch}{1.35}
\begin{longtable}{@{}p{0.7cm}p{3.5cm}p{5.6cm}p{5.4cm}@{}}
\toprule
& \textbf{Theorem} & \textbf{Hypotheses to verify} & \textbf{Gives you} \\
\midrule
\endhead
T1 & Neyman--Fisher factorisation & $f=g(\theta,T)h(\mathbf x)$, $h$ free of $\theta$ &
$T$ sufficient (and conversely) \\
T2 & Rao--Blackwell & $S$ sufficient; $\E T^2<\infty$ &
$\E[T\mid S]$ has same bias, no more variance \\
T3 & Lehmann--Scheff\'e & $S$ \textbf{complete} sufficient; $h(S)$ unbiased &
$h(S)$ is \emph{the} UMVUE, a.s.\ unique \\
T4 & UMVUE uniqueness & a UMVUE exists & it is a.s.\ unique (no completeness needed) \\
T5 & Exponential-family completeness & natural parameter space contains an open interval
(rectangle if $k$-dim) & $\sum_iT(X_i)$ complete sufficient \\
T6 & Basu & $S$ complete sufficient; $V$ ancillary & $S\ind V$ \\
T7 & Cram\'er--Rao & regularity: support free of $\theta$, differentiation under $\int$ legal &
$\Var(T)\ge[\psi']^2/(nI)$ \\
T8 & CRLB attainment & score $=k(\theta)[T-\psi(\theta)]$ & equality, iff exp.\ family \\
T9 & Bayes estimate & squared error loss & posterior mean; $r=\E_m[\Var(\theta\mid X)]$ \\
T10 & Constant-risk minimax & $T_\pi$ Bayes for a \textbf{proper} $\pi$; risk constant &
$T_\pi$ minimax \\
T11 & Limit of Bayes rules & $\pi_k$ proper; $r(\pi_k)\to c=\sup_\theta R(\theta,T^*)$ &
$T^*$ minimax \\
T12 & Neyman--Pearson & both hypotheses \textbf{simple} & MP test is the LR test \\
T13 & Exp.\ family $\Rightarrow$ MLR & $c(\theta)$ monotone & MLR in $T(x)$ \\
T14 & Karlin--Rubin & MLR in $T$; \textbf{one-sided} $H_A$ & $\{T>c\}$ is UMP; power monotone \\
T15 & Consistency via Chebyshev & bias $\to0$ and $\Var\to0$ & $T_n\pto\psi(\theta)$ \\
T16 & MLE consistency & identifiability, common support, uniform LLN, unique root &
$\hat\theta_n\pto\theta_0$ \\
\bottomrule
\end{longtable}

\section{Estimation}

\begin{Thm}{Neyman--Fisher factorisation}
\hyp $X_1,\dots,X_n$ have joint pdf/pmf $f(\mathbf x\mid\theta)$.

\conc $T$ is sufficient for $\theta$ $\iff$
$f(\mathbf x\mid\theta)=g\big(\theta,T(\mathbf x)\big)\,h(\mathbf x)$ with $h$ free of $\theta$.

\uses Every sufficiency claim in the course. Also proves that the posterior, the MLE and every
generalized Bayes estimate depend on the data only through $T$.
\end{Thm}

\begin{Thm}{Rao--Blackwell}
\hyp $S$ sufficient; $T$ an estimator with $\E_\theta T^2<\infty$.

\conc $T^*=\E[T\mid S]$ is a statistic, $\E_\theta T^*=\E_\theta T$, and
$R(\theta,T^*)\le R(\theta,T)$ for all $\theta$, strictly unless $T$ is already a function of
$S$.

\prf Sufficiency makes $\E[T\mid S]$ free of $\theta$, hence a statistic. Tower property gives
equal means, hence equal bias. Then
$\Var(T)=\E[\Var(T\mid S)]+\Var(\E[T\mid S])\ge\Var(T^*)$, the first term being $\ge0$ and
zero iff $T$ is $S$-measurable. Equal bias plus smaller variance gives smaller MSE. $\square$

\uses Constructing UMVUEs by conditioning a crude unbiased estimator (typically an indicator).
\end{Thm}

\begin{Thm}{Lehmann--Scheff\'e}
\hyp $S$ \textbf{complete} and sufficient; $h(S)$ unbiased for $\psi(\theta)$.

\conc $h(S)$ is the UMVUE of $\psi(\theta)$, a.s.\ unique.

\prf For any unbiased $T$, $T^*=\E[T\mid S]$ is unbiased, a function of $S$, and has
$\Var(T^*)\le\Var(T)$ by T2. Both $T^*$ and $h(S)$ are unbiased functions of $S$, so
$\E_\theta[T^*-h(S)]=0$ for all $\theta$ and completeness forces $T^*=h(S)$ a.s. Hence
$\Var(h(S))\le\Var(T)$. $\square$

\uses The workhorse. \textbf{Completeness is the hypothesis people forget} --- without it the
conclusion is false, as $U[\theta,\theta+1]$ shows.
\end{Thm}

\begin{Thm}{Uniqueness of the UMVUE}
\hyp A UMVUE exists.

\conc It is a.s.\ unique. (No complete sufficient statistic required --- this is independent
of T3.)

\prf Let $T_1,T_2$ be UMVUEs, both with variance $\sigma^2$. Then $T=\tfrac12(T_1+T_2)$ is
unbiased, so
$\sigma^2\le\Var(T)=\tfrac14[\sigma^2+\sigma^2+2\rho\sigma^2]=\tfrac{\sigma^2}{2}(1+\rho)$,
giving $\rho\ge1$; but $\rho\le1$, so $\rho=1$. Perfect correlation forces $T_2=aT_1+b$ with
$a>0$, and matching means and variances gives $a=1$, $b=0$. $\square$
\end{Thm}

\begin{Thm}{Completeness of exponential families}
\hyp $f(x\mid\theta)=p(\theta)h(x)\exp\{c(\theta)T(x)\}$ with support free of $\theta$, and the
natural parameter space $\Gamma=\{c(\theta):\theta\in\Theta\}$ contains an interval with
non-empty interior. ($k$-parameter version: a $k$-dimensional rectangle with non-empty
interior.)

\conc $S=\sum_i T(X_i)$ is complete sufficient.

\uses Establishing completeness in one line for Bin, Poi, Exp, Normal, Gamma, power-series and
Multinomial. \textbf{The hypothesis bites} for $N(\theta,\theta^2)$, where the natural
parameters trace a curve, not a rectangle.
\end{Thm}

\begin{Thm}{Basu}
\hyp $S$ complete sufficient for $\theta$; $V$ ancillary (distribution free of $\theta$).

\conc $S\ind V$ under every $P_\theta$.

\prf Let $g(s)=\Prob(V\le v\mid S=s)$, free of $\theta$ by sufficiency. Then
$\E_\theta[g(S)]=\Prob_\theta(V\le v)=\Prob(V\le v)$, a constant free of $\theta$ by
ancillarity. So $\E_\theta\big[g(S)-\Prob(V\le v)\big]=0$ for all $\theta$, and completeness
gives $g(S)=\Prob(V\le v)$ a.s. That is exactly independence. $\square$

\uses Proving $\Xbar\ind S^2$; independence of $\sum X_i$ and scale-invariant ratios (HW1 Q6);
and the range identity $\E[R]=\E[V]\E[S]$. \textbf{Spotting ancillarity:} differences are
ancillary in location families, ratios in scale families.
\end{Thm}

\section{Information}

\begin{Thm}{Cram\'er--Rao lower bound}
\hyp Regularity: the support does not depend on $\theta$, and differentiation under the
integral sign is legal.

\conc For any $T$ with $\E_\theta T=\psi(\theta)$,
$\ \Var_\theta(T)\ \ge\ \dfrac{[\psi'(\theta)]^2}{n\,I(\theta)}$, where
$I(\theta)=\Var_\theta\big(\tfrac{\partial}{\partial\theta}\log f\big)
=-\E_\theta\big[\tfrac{\partial^2}{\partial\theta^2}\log f\big]$.

\uses A benchmark, and an easy way to \emph{prove} something is the UMVUE when the bound is
attained. \textbf{Meaningless for $U[0,\theta]$} --- the support moves, regularity fails, and
$\tfrac{n+1}{n}X_{(n)}$ has variance of order $\theta^2/n^2$, far below any naive bound.
\end{Thm}

\begin{Thm}{Attainment of the CRLB}
\conc Equality holds at every $\theta$ $\iff$ the score factorises as
$\frac{\partial}{\partial\theta}\log f(\mathbf x\mid\theta)=k(\theta)\big[T(\mathbf x)-\psi(\theta)\big]$,
i.e.\ iff the model is a one-parameter exponential family and $\psi$ is the mean of its natural
sufficient statistic.

\uses Deciding attainment questions instantly (HW1 Q4: yes, for $\E X$ in a power-series
family; Mock Q3: no, for $e^{-\lambda}$ in the Poisson).

\begin{confusion}{failing the CRLB does not disqualify a UMVUE}
The CRLB bounds \emph{all} unbiased estimators. If no estimator attains it, the best one still
sits strictly above. ``Its variance exceeds the CRLB, so it is not the UMVUE'' is the single
most common error on this topic.
\end{confusion}
\end{Thm}

\section{Bayes and minimax}

\begin{Thm}{Bayes estimate under squared error loss}
\hyp Squared error loss. (Absolute error gives the posterior median; 0--1 loss gives the mode.)

\conc $T_\pi(x)=\E[\theta\mid X=x]$, the posterior mean, and
$r(\pi,T_\pi)=\E_m\big[\Var(\theta\mid X)\big]$.

\prf Write $f(x\mid\theta)\pi(\theta)=m(x)\pi(\theta\mid x)$ and swap the order of integration:
$r(\pi,T)=\int m(x)\big[\int(T(x)-\theta)^2\pi(\theta\mid x)d\theta\big]dx$. For each fixed $x$
the inner bracket is minimised over the number $T(x)$ at the posterior mean, with minimum value
the posterior variance. $\square$

\uses Every Bayes computation, and --- via the risk formula --- every minimax certificate.
\end{Thm}

\begin{Thm}{Constant-risk minimax criterion}
\hyp $T_\pi$ Bayes for a \textbf{proper} prior $\pi$; $R(\theta,T_\pi)=c$ constant in $\theta$.

\conc $T_\pi$ is minimax; $\pi$ is least favourable.

\prf $\sup_\theta R(\theta,T)\ge r(\pi,T)\ge r(\pi,T_\pi)=c=\sup_\theta R(\theta,T_\pi)$. $\square$

\uses The Binomial minimax estimator. \textbf{Properness is essential} --- under an improper
prior the Bayes risk is infinite and the middle step collapses.
\end{Thm}

\begin{Thm}{Limit of Bayes rules}
\hyp $\{\pi_k\}$ proper priors with $r(\pi_k,T_{\pi_k})\to c$; $\sup_\theta R(\theta,T^*)=c$.

\conc $T^*$ is minimax; $\{\pi_k\}$ is a least favourable sequence.

\prf $\sup_\theta R(\theta,T)\ge r(\pi_k,T)\ge r(\pi_k,T_{\pi_k})$ for every $k$; let
$k\to\infty$. $\square$

\uses $X$ minimax for $N(\theta,1)$ (HW2 Q1); and, run in reverse, proving the minimax value is
$\infty$ for the Poisson (HW2 Q2) and for $U[a,\theta]$.
\end{Thm}

\section{Testing}

\begin{Thm}{Neyman--Pearson lemma}
\hyp Both hypotheses \textbf{simple}: $H_0:X\sim f_0$ vs $H_A:X\sim f_1$.

\conc There exist $k\ge0,\gamma\in[0,1]$ with
$\phi^*=1,\gamma,0$ according as $f_1\gtrless kf_0$, having $\E_0\phi^*=\alpha$; any such
$\phi^*$ is most powerful of its size; and any MP level-$\alpha$ test has this form a.e.\ off
$\{f_1=kf_0\}$.

\prf On $\{f_1>kf_0\}$, $\phi^*=1\ge\phi$; on $\{f_1<kf_0\}$, $\phi^*=0\le\phi$; on
$\{f_1=kf_0\}$ the second factor vanishes. So $(\phi^*-\phi)(f_1-kf_0)\ge0$ pointwise.
Integrating,
$0\le(\E_1\phi^*-\E_1\phi)-k(\E_0\phi^*-\E_0\phi)\le\E_1\phi^*-\E_1\phi$. $\square$

\uses Every MP test. Remember the region is one-sided \emph{in $\Lambda$}, not necessarily in
$x$ --- see the Laplace-vs-normal annulus (HW2 Q5) and the non-monotone ordering in HW3 Q3.
\end{Thm}

\begin{Thm}{Exponential family $\Rightarrow$ MLR}
\hyp $f(x\mid\theta)=p(\theta)h(x)e^{c(\theta)T(x)}$ with $c$ non-decreasing.

\conc The family has MLR in $T(x)$.

\prf $\dfrac{f_{\theta_1}}{f_{\theta_0}}=\dfrac{p(\theta_1)}{p(\theta_0)}
\exp\{[c(\theta_1)-c(\theta_0)]T(x)\}$, non-decreasing in $T(x)$ since the exponent
coefficient is $\ge0$. $\square$

\uses Establishing MLR for Bin, Poi, Exp, $N(\theta,1)$, $N(0,\sigma^2)$ (in $x^2$).
$U[0,\theta]$ has MLR in $X_{(n)}$ but must be checked \emph{directly} --- it is not an
exponential family. The \textbf{Cauchy has no MLR}: the ratio of two quadratics tends to 1 at
both ends, so it cannot be monotone.
\end{Thm}

\begin{Thm}{Karlin--Rubin}
\hyp MLR in $T$; hypotheses \textbf{one-sided}, $H_0:\theta\le\theta_0$ vs $H_A:\theta>\theta_0$.

\conc $\phi^*=1,\gamma,0$ according as $T\gtrless c$, with $\E_{\theta_0}\phi^*=\alpha$, is UMP
of level $\alpha$; and $\beta_{\phi^*}(\theta)$ is non-decreasing, so the size is attained at
$\theta_0$.

\prf sketch: By NP the MP test against any fixed $\theta_1>\theta_0$ rejects for large
$f_{\theta_1}/f_{\theta_0}$, hence for large $T$ by MLR; the cutoff is determined by
$\Prob_{\theta_0}(T>c)=\alpha$, which does not involve $\theta_1$. So the same test is MP
against every $\theta_1>\theta_0$. Monotone power upgrades the null from $=\theta_0$ to
$\le\theta_0$ at no cost in size. $\square$

\uses All one-sided UMP tests (HW3 Q1, Q2; Mock Q8). \textbf{Randomise if $T$ is discrete.}

\begin{confusion}{two-sided alternatives}
No UMP exists in a regular MLR family: the best test against $\theta_1>\theta_0$ rejects large
$T$, against $\theta_1<\theta_0$ rejects small $T$. \textbf{Exception:} if the support moves
with $\theta$, the impossible-under-$H_0$ region is free, and a two-sided UMP test can exist ---
$U[0,\theta]$ with $H_0:\theta=1$ (HW3 Q4).
\end{confusion}
\end{Thm}

\section{Asymptotics}

\begin{Thm}{Consistency via Chebyshev}
\hyp $\E_\theta T_n\to\psi(\theta)$ and $\Var_\theta(T_n)\to0$.

\conc $T_n\pto\psi(\theta)$.

\prf $\Prob(|T_n-\psi|>\varepsilon)\le\varepsilon^{-2}\E(T_n-\psi)^2
=\varepsilon^{-2}[\text{bias}^2+\Var]\to0$. $\square$

\uses Almost every consistency question. \textbf{Sufficient, not necessary} --- $X_{(n)}$ for
$U[0,\theta]$ is biased at every $n$ and consistent; and an estimator with no moments at all
can be consistent.
\end{Thm}

\begin{Thm}{Consistency of the MLE (Wald)}
\hyp Identifiability; support free of $\theta$; $\Theta$ compact (or the maximiser eventually
in a compact set); a uniform law of large numbers for $\log f$.

\conc $\hat\theta_n\pto\theta_0$.

\prf sketch: By Jensen,
$M(\theta)-M(\theta_0)=\E_{\theta_0}\log\frac{f(X\mid\theta)}{f(X\mid\theta_0)}
\le\log\int f(x\mid\theta)dx=0$, with equality only at $\theta_0$ by identifiability --- so the
population criterion $M$ has a strict maximum at $\theta_0$. The sample criterion
$M_n=\frac1n\sum\log f(X_i\mid\theta)$ converges to $M$ uniformly, so its maximiser converges
to $\theta_0$. $\square$

\uses \textbf{Regularity is sufficient, not necessary}: $U[0,\theta]$ violates the
common-support hypothesis yet $X_{(n)}$ is consistent (and $n$-consistent, not merely
$\sqrt n$-consistent).
\end{Thm}

\begin{Thm}{Consistency of UMVUEs}
\hyp $T_n$ is the UMVUE of $\psi(\theta)$ for each $n$, and \emph{some} unbiased sequence
$\tilde T_n$ has $\Var(\tilde T_n)\to0$.

\conc $T_n\pto\psi(\theta)$.

\prf $0\le\Var(T_n)\le\Var(\tilde T_n)\to0$ by minimality; $T_n$ is unbiased, so Chebyshev
applies. $\square$
\end{Thm}

\section{Two results that are not called theorems but behave like them}

\textbf{Estimability in the Binomial.} $\psi(\theta)$ has an unbiased estimator based on
$X\sim\mathrm{Bin}(m,\theta)$ \emph{iff} $\psi$ is a polynomial of degree $\le m$. ($\Rightarrow$:
$\E_\theta T$ is such a polynomial, and two polynomials agreeing on an interval are identical.
$\Leftarrow$: $X^{\underline k}/m^{\underline k}$ is unbiased for $\theta^k$.)

\textbf{Estimability in the Poisson.} $\psi$ is estimable \emph{iff} $\psi$ is analytic, and
then the unbiased estimator is unique --- hence automatically the UMVUE. The same
uniqueness-by-transform argument runs through the Laplace transform for the exponential family.

\begin{idea}{The ``uniqueness $\Rightarrow$ UMVUE'' shortcut}
Whenever a series- or transform-matching argument produces the \emph{unique} unbiased
estimator, it is the UMVUE with nothing further to prove: there is no competitor. Worth one
sentence whenever it applies.
\end{idea}

\vfill
\begin{center}
\small\itshape
Companions: \texttt{PI\_Bayes\_Crash\_Course.pdf}, \texttt{PI\_Testing\_Crash\_Course.pdf},
\texttt{PI\_Exam\_Revision.pdf}, \texttt{PI\_Class\_Exercises\_Solved.pdf},
\texttt{PI\_Assignment\_Solutions.pdf}, \texttt{PI\_Mock\_Exam.pdf}.
\end{center}

\end{document}
